\clearpage
\section*{Дополнительные данные и проверки устойчивости}
% Generated by scripts/write_supplementary_numbers.py; do not hand-edit.
\newcommand{\SN}[1]{\csname supp@#1\endcsname}
\expandafter\def\csname supp@followup.total.M\endcsname{3.00}
\expandafter\def\csname supp@followup.total.F\endcsname{2.07}
\expandafter\def\csname supp@followup.urban.M\endcsname{-0.32}
\expandafter\def\csname supp@followup.urban.F\endcsname{-8.10}
\expandafter\def\csname supp@followup.rural.M\endcsname{3.49}
\expandafter\def\csname supp@followup.rural.F\endcsname{3.78}
\expandafter\def\csname supp@urban.b\endcsname{-7.66}
\expandafter\def\csname supp@urban.ci_low_cr2\endcsname{-12.48}
\expandafter\def\csname supp@urban.ci_high_cr2\endcsname{-2.83}
\expandafter\def\csname supp@urban.wcr_p\endcsname{0.0625}
\expandafter\def\csname supp@policy.18-37.female.b\endcsname{20.354}
\expandafter\def\csname supp@policy.18-37.female.ci_low_cr2\endcsname{16.679}
\expandafter\def\csname supp@policy.18-37.female.ci_high_cr2\endcsname{24.028}
\expandafter\def\csname supp@policy.18-37.DDD.b\endcsname{14.379}
\expandafter\def\csname supp@policy.18-37.DDD.ci_low_cr2\endcsname{9.739}
\expandafter\def\csname supp@policy.18-37.DDD.ci_high_cr2\endcsname{19.020}
\expandafter\def\csname supp@policy.18-32.female.b\endcsname{20.312}
\expandafter\def\csname supp@policy.18-32.female.ci_low_cr2\endcsname{16.606}
\expandafter\def\csname supp@policy.18-32.female.ci_high_cr2\endcsname{24.019}
\expandafter\def\csname supp@policy.18-32.DDD.b\endcsname{14.377}
\expandafter\def\csname supp@policy.18-32.DDD.ci_low_cr2\endcsname{9.745}
\expandafter\def\csname supp@policy.18-32.DDD.ci_high_cr2\endcsname{19.009}

В обновлённую репликационную базу включены региональные таблицы переписи
1959~г. и обследование учреждений ликвидации безграмотности на 1 января
1924~г. В последнем визуально проверены 30 строк по 15 территориям;
14 территорий соединены с проверенной когортой 18--22 лет в 1926~г.
Показатели относятся к учреждениям и учащимся на дату обследования,
а не к уникальным выпускникам или охвату определённой возрастной группы.
Неодинаковая полнота обследования, изменения границ и эндогенность
размещения учреждений ограничивают интерпретацию. Во всех 24 разведочных
моделях гендерного состава и масштаба сети интервалы CR2 включают ноль;
надёжная причинная декомпозиция каналов на этих данных не получена.

В плане на 1931~г. возрастной контингент сельского ликбеза определён
как 16--40 лет, с отдельными исключениями для старших возрастов.
Исключение группы 33--37 лет в 1926~г., пересекающей границу 40 лет
к 1931~г., исключает её из промежуточной группы, сохраняя молодую группу 18--32 лет, и меняет основной женский контраст с
\SN{policy.baseline.female.b} до \SN{policy.omit_33_37.female.b} п.п.
Интервал CR2 после исключения промежуточной когорты равен
[\SN{policy.omit_33_37.female.ci_low_cr2}; \SN{policy.omit_33_37.female.ci_high_cr2}].
Это проверка одной возрастной границы, а не доказательство отсутствия
обучения в старшей контрольной группе.

Региональные сопоставления 1939--1959~гг. выявляют неоднородность
динамики между городом и селом. Эти агрегаты не прослеживают отдельных
людей и не позволяют ограничить величину ошибки переписи 1939~г. без
дополнительных предпосылок об обучении, миграции и отборе. Полные источники,
оценки, альтернативные выборки и ограничения приведены в двух папках
\foreignlanguage{english}{\texttt{supplementary}} репликационного проекта.
